\chapter*{\MakeUppercase{Resumen}}
%\chapter*{Resumen}
\thispagestyle{fancy}
%\addcontentsline{toc}{chapter}{\MakeUppercase{Resumen}}
%\addcontentsline{toc}{chapter}{Resumen}
\noindent El presente trabajo tiene como objetivo encontrar una aproximación numérica de la evolución de la frontera de un tumor avascular, para ello se emplea el Método de Elementos Finitos en la solución de problemas elípticos que determinan la presión y la concentración de nutrientes al interior del tumor; la frontera del tumor está representada por una curva de nivel de una función continua $\phi$, y su evolución está determinada por el gradiente de la presión. Finalmente se realizan simulaciones para diferentes fronteras iniciales del tumor.

\noindent Palabras clave -- Crecimiento tumoral, método de conjunto de nivel, método de elementos  finitos, frontera libre.
\chapter*{\MakeUppercase{Abstract}}
%\chapter*{Abstract}
\thispagestyle{fancy}
%\addcontentsline{toc}{chapter}{\MakeUppercase{Abstract}}
%\addcontentsline{toc}{chapter}{Abstract}
\noindent The objective of this work is to find a numerical approximation of the evolution of the boundary of an avascular tumor, using the Finite Element Method in the solution of elliptic problems that determine the pressure and the concentration of nutrients inside the tumor; the tumor boundary is represented by a level curve of a continuous function $\phi$ and its evolution is determined by the pressure gradient. Finally, simulations are performed for different initial tumor boundaries.

\noindent Keyword --  Tumor growth, level set method, finite element method, free boundary.